% Page furniture: chapter openers, headings, running heads, tables, code and callouts. Colours come from
% tokens/colors.tex, which is generated from brand.yml.
\usepackage{titlesec}
\usepackage{fancyhdr}
\usepackage{graphicx}
\usepackage{tikz}
\usepackage{eso-pic}

% A document sets its own part or section tagline with \renewcommand in
% header-includes, which lands in the preamble. cover.tex runs from
% include-before-body, i.e. after \begin{document}, so a \providecommand there
% would be too late and the \renewcommand fails with "Command \booktagline
% undefined". It has to be declared from an include-in-header file, which this
% is.
\providecommand{\booktagline}{}

% The cover and the back page both read these, and the cover runs first, so all
% four have to be declared here rather than in either of those files. This is
% the same trap \booktagline sits in above: include-before-body and
% include-after-body both land after the preamble, so a \providecommand in
% either is too late for anything earlier to use, and for anything a document
% sets in header-includes to override.
\providecommand{\bookweb}{}
\providecommand{\bookrepo}{}
\providecommand{\bookpublisher}{}
% The running head prints the book's title; the class clears \@title once the title
% page is set, so it is kept here and set from header-includes.
\providecommand{\booktitle}{}

% Quarto emits \maketitle even under a custom cover, and the class clears
% \@title, \@author and \@date once it has typeset them. The cover runs from
% include-before-body, which lands *after* that, so it would find all three
% empty. Neutralising \maketitle removes the duplicate title page and leaves the
% macros intact for the cover to read. Deferred to \AtBeginDocument because
% Quarto's template redefines \maketitle after this file is included.
\makeatletter
\AtBeginDocument{\let\maketitle\relax}
\makeatother

% ---------------------------------------------------------------------------
% Chapter openers
% ---------------------------------------------------------------------------
%
% A large muted numeral above the title with a blue rule beneath, rather
% than KOMA's default. `display` is used so the number and title are separate
% blocks; the numeral is set in the cover weight because at 42pt the heading
% weight looks thin.
\titleformat{\chapter}[display]
  {\bkHeadingFont\color{knInk}}
  {\bkCoverFont\fontsize{42}{46}\selectfont\color{knSlateLight}\thechapter}
  {6pt}
  {\Huge\bfseries}
  [{\vspace{6pt}\color{knBlue}\titlerule[2pt]}]
\titlespacing*{\chapter}{0pt}{10pt}{22pt}

\titleformat{\section}
  {\bkHeadingFont\Large\bfseries\color{knInk}}{\thesection}{0.6em}{}
  [{\vspace{2pt}\color{knRule}\titlerule[1.2pt]}]
\titleformat{\subsection}
  {\bkHeadingFont\large\bfseries\color{knInk}}{\thesubsection}{0.6em}{}
% Without this a `###` falls back to the class default, which in the body face
% is indistinguishable from a bold paragraph.
\titleformat{\subsubsection}
  {\bkHeadingFont\normalsize\bfseries\color{knInk}}{\thesubsubsection}{0.6em}{}
\titlespacing*{\section}{0pt}{16pt}{6pt}
\titlespacing*{\subsection}{0pt}{12pt}{4pt}
\titlespacing*{\subsubsection}{0pt}{10pt}{3pt}

% ---------------------------------------------------------------------------
% Running heads
% ---------------------------------------------------------------------------
\pagestyle{fancy}
\fancyhf{}
\renewcommand{\headrulewidth}{0.4pt}
\renewcommand{\footrulewidth}{0pt}
% \MakeUppercase is deliberately absent: Jost's capitals are wide, and a long
% chapter title set in caps wraps the running head onto two lines.
\fancyhead[L]{\bkHeadingFont\footnotesize\color{knSlate}\leftmark}
\fancyhead[R]{\bkHeadingFont\footnotesize\color{knSlate}\booktitle}
\fancyfoot[C]{\bkHeadingFont\footnotesize\color{knSlate}\thepage}

% KOMA sets \chaptermark with the number prefixed; the running head reads better
% with the title alone, since the number is already on the opening page.
\renewcommand{\chaptermark}[1]{\markboth{#1}{}}

% The cover has no running head and no page number.
\fancypagestyle{bkcover}{\fancyhf{}\renewcommand{\headrulewidth}{0pt}}

% A chapter's own opening page uses `plain` in KOMA, which would print a page
% number under a design that has none elsewhere on that page. Matched to the
% rest instead.
\fancypagestyle{plain}{%
  \fancyhf{}%
  \renewcommand{\headrulewidth}{0pt}%
  \fancyfoot[C]{\bkHeadingFont\footnotesize\color{knSlate}\thepage}%
}

% ---------------------------------------------------------------------------
% Tables
% ---------------------------------------------------------------------------
%
% Blue rules, as the web edition gives a table a blue top border. A filled header row is not available on
% this path: \rowcolor needs xcolor's `table` option, Quarto loads xcolor before
% include-in-header runs, and re-loading with the option is an option clash.
\usepackage{booktabs}
\usepackage{colortbl}
\arrayrulecolor{knBlue}
\setlength{\heavyrulewidth}{1.2pt}
\setlength{\lightrulewidth}{0.8pt}
\renewcommand{\arraystretch}{1.15}

% ---------------------------------------------------------------------------
% Code blocks
% ---------------------------------------------------------------------------
%
% A pale surface with a blue left edge, as the web edition, so a listing reads as a distinct
% object without a heavy box around it. Quarto defines Shaded around every
% highlighted block, so redefining that environment reaches all of them.
\usepackage{tcolorbox}
\tcbuselibrary{skins, breakable}

% `Shaded` wraps every highlighted listing, and Quarto defines it only when a
% document actually contains one, at a point that varies with the format. A bare
% \renewenvironment therefore works for a book render and dies with "Environment
% Shaded undefined" for a single document. Deferring to \AtBeginDocument and
% choosing between \newenvironment and \renewenvironment covers both.
\makeatletter
\AtBeginDocument{%
  \@ifundefined{Shaded}{\newenvironment}{\renewenvironment}{Shaded}%
    {\begin{tcolorbox}[
       breakable, enhanced, sharp corners,
       colback=knPanel, colframe=knPanel,
       borderline west={2.2pt}{0pt}{knBlue},
       boxrule=0pt, left=8pt, right=8pt, top=6pt, bottom=6pt,
     ]}%
    {\end{tcolorbox}}%
}
\makeatother

% Quarto callouts (note, tip, warning, important, caution) are tcolorbox environments; they take the brand colours.
\AtBeginDocument{%
  \colorlet{quarto-callout-note-color}{knBlue}\colorlet{quarto-callout-note-color-frame}{knBlue}%
  \colorlet{quarto-callout-tip-color}{knSuccess}\colorlet{quarto-callout-tip-color-frame}{knSuccess}%
  \colorlet{quarto-callout-warning-color}{knOrange}\colorlet{quarto-callout-warning-color-frame}{knOrange}%
  \colorlet{quarto-callout-important-color}{knDanger}\colorlet{quarto-callout-important-color-frame}{knDanger}%
  \colorlet{quarto-callout-caution-color}{knOrange}\colorlet{quarto-callout-caution-color-frame}{knOrange}%
}

% Links, citations and contents entries in the brand blue, not hyperref's default red.
\AtBeginDocument{\hypersetup{linkcolor=knBlue, citecolor=knBlue, urlcolor=knBlue, filecolor=knBlue}}

% The contents list in ink, not in the link colour: its entries are links, and blue entries read as a page of hyperlinks.
\usepackage{etoolbox}
\pretocmd{\tableofcontents}{\hypersetup{linkcolor=knInk}}{}{}
\apptocmd{\tableofcontents}{\hypersetup{linkcolor=knBlue}}{}{}

% A box of a book (remark, pitfall, key concept ...): a tinted block with a coloured left edge and a title in the heading
% font. The title is set inside the box, as on the web, where it sits on the same tint as the text.
\usepackage{tcolorbox}
\tcbuselibrary{skins, breakable}
\newtcolorbox{knbox}[4]{%
  breakable, enhanced, sharp corners,
  colback=#2, colframe=#2, boxrule=0pt,
  borderline west={3.5pt}{0pt}{#1},
  left=11pt, right=11pt, top=9pt, bottom=9pt,
  before skip=10pt, after skip=10pt,
  before upper={{\bkHeadingFont\bfseries\color{#3}#4\par}\vspace{5pt}},
}
